% =====================================================================
%  Thème 5 — CORRIGÉ — Niveau MOYEN
% =====================================================================
\documentclass[11pt,a4paper]{article}
\usepackage[utf8]{inputenc}
\usepackage[T1]{fontenc}
\usepackage{lmodern}
\usepackage[french]{babel}
\usepackage{amsmath,amssymb,mathtools}
\usepackage{icomma}
\usepackage{xcolor}
\usepackage[most]{tcolorbox}
\usepackage{enumitem}
\usepackage{array}
\usepackage{booktabs}
\usepackage{fancyhdr}
\usepackage[hidelinks]{hyperref}
\usepackage[a4paper,margin=2.2cm,headheight=26pt,headsep=10pt]{geometry}

\definecolor{primary}{RGB}{223,87,99}
\definecolor{primarydark}{RGB}{150,42,55}
\definecolor{excol}{RGB}{0,135,90}

\tcbset{boxbase/.style={enhanced,breakable,boxrule=0.9pt,arc=2.5pt,
   left=3mm,right=3mm,top=2.3mm,bottom=2.3mm,
   fonttitle=\bfseries,coltitle=white,toptitle=0.6mm,bottomtitle=0.6mm}}
\newtcolorbox[auto counter]{correction}[1][]{boxbase,colback=excol!5,colframe=excol,
  title={\textsc{Corrigé}~\thetcbcounter\ifx\relax#1\relax\relax\else\textmd{ --- \itshape #1}\fi}}

\pagestyle{fancy}
\fancyhf{}
\fancyhead[L]{\small\textcolor{primarydark}{\textbf{Fiche 6} — Les limites}}
\fancyhead[R]{\small\textcolor{primarydark}{\textsc{Thème 5 — Corrigé Moyen}}}
\fancyfoot[C]{\small\thepage}
\renewcommand{\headrulewidth}{0.8pt}
\renewcommand{\headrule}{\hbox to\headwidth{\color{primary}\leaders\hrule height \headrulewidth\hfill}}

\newcommand{\R}{\mathbb{R}}

\begin{document}

\begin{center}
{\Large\bfseries\color{primarydark}Thème 5 — L'Hospital \& limites trigonométriques}\\[2pt]
{\large\color{excol}Corrigé — Niveau Moyen}
\end{center}
\vspace{6pt}

\begin{correction}[l'Hospital avec le cosinus]
\begin{enumerate}[label=\alph*),leftmargin=2em,itemsep=3pt]
  \item En $\frac{\pi}{2}$ : numérateur $\cos\frac{\pi}{2}=0$, dénominateur $\frac{\pi}{2}-\frac{\pi}{2}=0$ : forme $\frac00$.
  \item $\displaystyle\lim_{x\to \frac{\pi}{2}}\frac{\cos x}{x-\frac{\pi}{2}}=\lim_{x\to\frac{\pi}{2}}\frac{-\sin x}{1}=-\sin\frac{\pi}{2}=-1$.
        \emph{(C'est aussi le nombre dérivé de $\cos$ en $\frac{\pi}{2}$.)}
\end{enumerate}
\end{correction}

\begin{correction}[combinaison de limites remarquables]
\[ \frac{\sin 2x}{\tan x}=\frac{\sin 2x}{2x}\times\frac{x}{\tan x}\times\frac{2x}{x}
   =\underbrace{\frac{\sin 2x}{2x}}_{\to 1}\times\underbrace{\frac{x}{\tan x}}_{\to 1}\times 2
   \xrightarrow[x\to0]{}1\times 1\times 2=2. \]
\end{correction}

\begin{correction}[transformer $0\times\infty$]
\begin{enumerate}[label=\alph*),leftmargin=2em,itemsep=3pt]
  \item En $0^{+}$ : $x\to 0$ et $\ln x\to-\infty$ : forme $0\times(-\infty)$, indéterminée (le produit
        dépend de la « vitesse » relative des deux facteurs).
  \item $\displaystyle x\ln x=\frac{\ln x}{1/x}$ (forme $\frac{-\infty}{+\infty}$). Par l'Hospital :
        \[ \lim_{x\to 0^{+}}\frac{\ln x}{1/x}=\lim_{x\to 0^{+}}\frac{1/x}{-1/x^2}=\lim_{x\to 0^{+}}(-x)=0. \]
\end{enumerate}
\end{correction}

\begin{correction}[l'Hospital deux fois]
\begin{enumerate}[label=\alph*),leftmargin=2em,itemsep=3pt]
  \item Forme $\frac{1-1}{0}=\frac00$. Une fois : $\displaystyle\lim_{x\to 0}\frac{1-\cos x}{x^2}=\lim_{x\to 0}\frac{\sin x}{2x}$.
  \item C'est encore $\frac00$. Seconde fois : $\displaystyle\lim_{x\to 0}\frac{\sin x}{2x}=\lim_{x\to 0}\frac{\cos x}{2}=\frac{\cos 0}{2}=\frac12$.
        \emph{(On pouvait aussi conclure dès l'étape 1 : $\frac{\sin x}{2x}=\frac12\cdot\frac{\sin x}{x}\to\frac12$.)}
\end{enumerate}
\end{correction}

\begin{correction}[la même expression, deux points]
$q(x)=\dfrac{e^x-1}{x}$.
\begin{enumerate}[label=\alph*),leftmargin=2em,itemsep=3pt]
  \item $\frac00$ : $\displaystyle\lim_{x\to 0}\frac{e^x-1}{x}=\lim_{x\to 0}\frac{e^x}{1}=1$.
  \item $\frac{\infty}{\infty}$ : $\displaystyle\lim_{x\to+\infty}\frac{e^x-1}{x}=\lim_{x\to+\infty}\frac{e^x}{1}=+\infty$.
  \item Près de $0$, $e^x-1\approx x$ (le rapport tend vers $1$) : les deux croissent « à la même vitesse ».
        Mais à l'infini, $e^x$ l'emporte largement sur $x$ (le rapport explose) : \textbf{l'exponentielle croît
        plus vite que toute puissance de $x$}.
\end{enumerate}
\end{correction}

\end{document}
